\section{Acceso al Panel Gerencial y Control de Cuentas}

El Panel de Administración es un portal de acceso restringido diseñado exclusivamente para ti y el grupo de supervisores.

\subsection{El Portal Secreto}

\begin{enumerate}
    \item En tu navegador de internet, escribe la dirección directa al panel gerencial (generalmente añadiendo \texttt{/admin} al final de la dirección de la clínica).
    \item Aparecerá una pantalla de inicio de sesión especial color azul.
    \item \textbf{Nota Importante:} Las credenciales ejecutivas (tu usuario y contraseña con nivel de gerente) fueron creadas y te deben ser entregadas personalmente por el Equipo de Sistemas (TI) al momento de inaugurar e instalar la plataforma.
    \item Una vez que las ingreses, verás el menú con todas las categorías de información de la clínica.
\end{enumerate}

\begin{figure}[H]
    \centering
    \includegraphics[width=0.6\linewidth]{screenshots/admin_login.png}
    \caption{Pantalla de acceso ejecutivo.}
\end{figure}

\subsection{Encontrar Usuarios y Cambiar Contraseñas}

A este sistema entran desde pacientes hasta doctores radiólogos. Todos ellos tienen una cuenta base llamada "Usuario". Si un paciente llama por teléfono diciendo que olvidó su clave, tú puedes ayudarlo:

\begin{enumerate}
    \item En el menú principal, busca la sección \textbf{Authentication and Authorization} (Autenticación) y da clic en \textbf{Users} (Usuarios).
    \item Verás una lista con todos los correos electrónicos registrados. Si la lista es muy larga, usa la barra de búsqueda en la parte superior.
    \item Da clic sobre el correo electrónico de la persona a la que deseas ayudar.
\end{enumerate}

\textbf{Para Cambiarle la Contraseña manualmente:}
Mira fijamente al principio de su página. Verás un texto azul que dice \textit{"this form"} (este formulario). Haz clic ahí y el sistema te pedirá escribir una contraseña nueva dos veces. Guárdala y envíasela al individuo por un medio seguro.

\textbf{Definirles un Rol del menú:}
Si bajas un poco, verás un área llamada \textbf{Permissions} (Permisos). Del lado izquierdo están los roles posibles (\textit{Patients}, \textit{Assistants}, \textit{Doctors}, etc.). Puedes seleccionar uno, dar clic a la flechita y pasarlo a la columna de la derecha. Esto le indicará al sistema qué pantalla mostrarle la próxima vez que esa persona inicie sesión.
